\section{What source validation can identify}
\label{sec:selection-theory}

An adaptation rule can change the target decision function without changing
any source-validation prediction. This section gives an exact example and
then separates it from transformations for which the invariance does not hold.
Features are row vectors. All maps are fitted on source training and
unlabelled target adaptation data, never on target test data.

\subsection{A source-translation invariance}
\label{sec:translation-theory}

Consider a candidate offset $\delta$. Add it to every source-training and
source-validation vector, but leave target vectors unchanged. Mean matching
chooses $\delta=\mu_t-\mu_s$; the argument holds for any fixed offset.

\paragraph{Proposition 1: stationary-kernel validation.}
Let a kernel classifier use $k(x,y)=\psi(x-y)$. Suppose fitting and tie-breaking
depend on the source features only through the training Gram matrix, with
identical labels, hyperparameters and random state. Translating all source
vectors by $\delta$ leaves every source-validation score unchanged. Its target
decision function satisfies
\begin{equation}
 f_\delta(x)=f_0(x-\delta).
 \label{eq:translation-kernel}
\end{equation}
This statement holds for each candidate in a fixed hyperparameter search, not
only for the selected candidate.

\paragraph{Proof.}
Training entries obey $k(x_i+\delta,x_j+\delta)=k(x_i,x_j)$. Validation entries
obey the same identity because both arguments are translated. The fitted
coefficients, intercept and validation predictions are therefore unchanged
under the stated fitting rule. For an unshifted target vector,
$k(x_i+\delta,x)=k(x_i,x-\delta)$. Substitution into the fitted kernel expansion
gives Equation~\ref{eq:translation-kernel}. The argument applies separately
to each binary decision function in a multiclass SVM.

\paragraph{Consequence for selection.}
Even an arbitrarily large labelled source-validation set cannot distinguish
these offsets using the classifier's predictions on that set. A deterministic
tie-break can choose an offset, but its validation score supplies no evidence
that this offset has the best target risk. Target predictions can change when
$x-\delta$ crosses a decision boundary. This is a limitation of this criterion,
not an impossibility result for selectors that also use unlabelled target
features, additional assumptions, or target supervision.

\paragraph{Linear objectives and numerical limits.}
The same phenomenon occurs for an affine head $xW+q$ trained with a loss on
source predictions and a penalty on $W$ alone. Replacing $q$ by $q-\delta W$
preserves every translated-source score and the penalty. Corresponding optima
therefore have the same source risk but can have different target scores.
This is an objective identity, not a guarantee about finite optimisation.
Penalising the intercept breaks it. The implemented \texttt{LinearSVC} uses a
penalised synthetic intercept; neural initialisation, weight decay and early
stopping also prevent an unrestricted invariance claim. Logistic regression
with an unpenalised intercept satisfies the ideal objective assumptions, but
its finite solver can return different numerical solutions. We report that
contrast rather than treating it as an exception to be discarded.

\subsection{Separating transport from classifier coordinates}
\label{sec:affine-factorisation}

The translation result does not imply that validation is invariant to
standardisation or CORAL. To state the distinction, write the two domain maps
as $T_s(x)=xA+a$ and $T_t(x)=xB+b$, where $B$ is invertible. Define
\begin{equation}
 R=AB^{-1},\qquad r=(a-b)B^{-1}.
 \label{eq:relative-transport}
\end{equation}
Then $T_s(x)=(xR+r)B+b$. Thus $xR+r$ is the source transport expressed in
original target coordinates, followed by a common coordinate map $zB+b$.
The source matrix $A$ need not be invertible. This covers source-only affine
maps, whose target map is the identity.

\paragraph{Lemma 2: the induced geometries.}
For a head $zW+q$, let $V=BW$, $u=bW+q$, and $G=BB^\top$. Source scores become
$(xR+r)V+u$ and target scores become $xV+u$. A Euclidean coefficient penalty
and an RBF distance in the transformed coordinates become
\begin{align}
 \lVert W\rVert_F^2 &= \operatorname{tr}(V^\top G^{-1}V),
 \label{eq:penalty-pullback}\\*
 \lVert(x-y)B\rVert^2 &= (x-y)G(x-y)^\top.
 \label{eq:kernel-pullback}
\end{align}
The identities follow by substituting $W=B^{-1}V$ and expanding the squared
norm. The common map therefore induces inverse metrics for the coefficient
penalty and the feature distance. They are not additional alignment methods.

These elementary reparameterisations are not presented as new linear algebra.
Feature/weight equivalence is already part of the CORAL account
\cite{sun2016coral}. Their use here is to specify what the experimental
comparison holds fixed. Equal trial counts and equal scalar regularisation
ranges do not generally equalise an anisotropic penalty matrix. Nor does a
scalar RBF bandwidth search remove arbitrary anisotropy. Equivalence of
objectives also does not imply identical optimisation trajectories.

\paragraph{Why z-score is not a moment-only control.}
Let $D_s,D_t$ contain the effective positive standard deviations used by the
implementation, including its constant-dimension replacements. Domain-wise
z-score has
\begin{equation}
 \begin{split}
 R&=D_s^{-1}D_t,\qquad r=\mu_t-\mu_sR,\\
 B&=D_t^{-1},\qquad b=-\mu_tD_t^{-1}.
 \end{split}
 \label{eq:zscore-factorisation}
\end{equation}
It is diagonal moment transport followed by common target standardisation.
In raw target coordinates its weight penalty uses $D_t^2$, while its distance
metric uses $D_t^{-2}$. A classifier with an isotropic penalty fitted directly
on diagonal moment-matched features is generally a different learner.
Consequently, the existing alignment comparison measures complete pipelines.
It does not isolate the causal contribution of matching additional moments.

\subsection{What a discrepancy comparison measures}
\label{sec:audit-measurement}

Equation~\ref{eq:kernel-pullback} also applies to an MMD kernel. For a change
of common coordinates, keeping Euclidean distance in the new coordinates
changes the kernel in the old coordinates. Conversely, if the intention is
only to express a fixed metric $H$ in coordinates $z'=zQ$, its representation
must change to $Q^{-1}HQ^{-\top}$. These are different measurement questions.
Forcing invariance to every invertible map is not a universal solution;
representation-comparison research shows that excessive invariance can discard
useful information \cite{kornblith2019similarity}.

We therefore audit three objects separately: relative transport $(R,r)$,
the classifier objective in those coordinates, and the diagnostic kernel.
The existing corpus scores cannot retrospectively separate their causal
effects. They can test the exact translation prediction, demonstrate the
failure of a particular source-validation criterion, and quantify sensitivity
to specified measurement protocols. No new alignment algorithm is required
for those questions.
